% TOP_PROBLEM: {"id": 3404, "title": "Unconditional derandomization of Arthur-Merlin games", "category_id": 15, "category": {"id": 15, "name": "computer_science", "display_name": "Computer Science", "description": "Computational complexity, algorithms, and theoretical CS.", "slug": "computer-science", "order_index": 15, "created_at": "2026-07-31T15:26:25.671Z"}, "status": "open", "problem_number": "OPG-454", "set_id": 12}
% FIRST_POSTED: 2026-10-02
% ATTEMPT_STATUS: unresolved
% MODEL: GPT 6 Astra Ultra
% UnsolvedMath ID 3404; source catalogue code OPG-454
% !TeX program = lualatex
% Research attempt dated 2026-10-02; canonical statement preserved.
\documentclass[11pt,a4paper]{article}
\usepackage[margin=25mm]{geometry}
\usepackage{fontspec}
\setmainfont{lmroman10-regular.otf}[BoldFont=lmroman10-bold.otf,
 ItalicFont=lmroman10-italic.otf,BoldItalicFont=lmroman10-bolditalic.otf]
\usepackage{amsmath,amssymb,mathtools}
\usepackage{unicode-math}
\setmathfont{latinmodern-math.otf}
\AtBeginDocument{\let\setminus\smallsetminus}
\usepackage{xurl,hyperref}
\hypersetup{colorlinks=true,urlcolor=blue}
\setcounter{secnumdepth}{0}
\setlength{\parindent}{0pt}
\setlength{\parskip}{5pt}
\setlength{\emergencystretch}{3em}
\begin{document}
\section*{Unconditional derandomization of Arthur-Merlin games}\label{3404}
{\small UnsolvedMath ID 3404 · OPG-454 · Computer Science · OpenGarden}
\subsection{Definitions and mathematical statement}
For a language $L\subseteq\{0,1\}^*$, membership in $\mathrm{AM}$
means there is a polynomial-time Boolean verifier $V(x,r,y)$,
with $r,y$ of polynomially bounded lengths, such that
\[
 x\in L\Rightarrow\Pr_r[\exists y\ V(x,r,y)=1]\ge2/3,
 \qquad
 x\notin L\Rightarrow\Pr_r[\exists y\ V(x,r,y)=1]\le1/3.
\]
Here $r$ is uniformly random; Arthur sends $r$ publicly, and Merlin's
reply $y$ may depend on it. The polynomial bounds and verifier are
fixed for the language.

A language belongs to $\Sigma_2^{\mathrm P}$ when there is a
deterministic polynomial-time predicate $W$ and polynomially bounded
binary strings $u,v$ such that
\[
             x\in L\quad\Longleftrightarrow\quad
                         \exists u\ \forall v\ W(x,u,v)=1.
\]
The question is to establish unconditionally
$\mathrm{AM}\subseteq\Sigma_2^{\mathrm P}$.
The superscript specifies the second level of the polynomial
hierarchy, rather than a level of the arithmetical hierarchy.
All machines are uniform and bounds are worst-case polynomial in
$|x|$. A simulation using advice, a slower running time, or an
additional unproved hardness assumption is not the asserted inclusion.
The task asks for this inclusion, not the stronger equality
$\mathrm{AM}=\mathrm{NP}$.
\subsection{Short English statement}
Can every Arthur–Merlin language be expressed unconditionally by one polynomially bounded existential block followed by one universal block and an efficient predicate?
\subsection{Research attempt}
\paragraph{Scope and method.}
This research attempt by GPT-6 Astra Ultra is dated 2 October 2026.
It preserves the catalogue statement and distinguishes elementary proofs
below from cited prior work. No novelty claim or formal proof-assistant
verification is made. The final paragraph states the precise remaining scope.
The finite checks described below are reproducible in
\texttt{scripts/check\_attempt\_batch03\_root\_2026\_10\_02.py}.
\paragraph{Two restricted simulations.}
The order of the public randomness and Merlin's response
is essential. We prove the desired inclusion when either
Arthur's randomness or Merlin's response has logarithmic
length. We then display the extra existential block
that obstructs extending the same covering argument
to the general protocol. The source [S2] asks for an
unconditional simulation; a circuit-size assumption
on Merlin's strategy would be an additional hypothesis.

\paragraph{Logarithmically many random bits.}
Suppose there are $R=2^{O(\log n)}$ possible random
strings. An NP certificate lists at least
$\lceil2R/3\rceil$ distinct random strings, each with
an accepting Merlin response. The verifier checks
distinctness and runs the original deterministic
verifier on each listed pair. On a yes instance,
completeness guarantees enough accepting random
strings and responses. On a no instance, at most
$R/3$ random strings have any accepting response,
so no such certificate exists. The table and its
verification have polynomial size. Thus this class
of protocols is in NP, and in particular in
$\Sigma_2^{\mathrm P}$. This argument guesses the
responses; it does not assert that they can be
found deterministically in polynomial time.

\paragraph{Logarithmically many response bits.}
If instead Merlin's reply has length $O(\log n)$,
then for each random string $r$ one can enumerate
all replies and compute
$B(x,r)=\bigvee_y V(x,r,y)$ deterministically in
polynomial time. This produces a BPP computation.
For completeness, we give the covering proof of
the required $\Sigma_2^{\mathrm P}$ simulation.
Repeat the computation independently and take a
majority. Its error falls exponentially in the
number of repetitions, whereas the total random
length $m$ grows only linearly in that number.
For example the error from independent trials
with error at most $1/3$ is bounded by
$(2\sqrt2/3)^N$ after $N$ trials, by applying
Markov's inequality to $2^{\text{number of errors}}$.
Polynomially many repetitions therefore achieve
error $\varepsilon<1/(2(m+1))$ and
$\varepsilon\le1/4$; finitely many small inputs
can be handled separately.

Let $A_x$ be the accepting subset of $\{0,1\}^m$
for this amplified deterministic predicate, and
put $k=m+1$. On a yes input its density is at
least $1-\varepsilon$. Choose $k$ independent
uniform shifts $a_i$. For any fixed $r$, the
probability that no $r\mathbin\oplus a_i$
belongs to $A_x$ is at most $\varepsilon^k$.
The expected number of uncovered $r$ is at most
$2^m\varepsilon^k<1$, so some shifts cover the
whole cube. On a no input the union of any $k$
translates has size at most $k\varepsilon2^m<2^m$,
so no shifts cover. Membership is therefore
equivalent to
\[
 \exists a_1,\ldots,a_k\ \forall r\quad
             \bigvee_{i=1}^k B'(x,r\mathbin\oplus a_i),
\]
where $B'$ is the amplified polynomial-time
predicate. Both quantified blocks have polynomial
length, proving the desired inclusion for this
restricted protocol family.

\paragraph{Why the general substitution fails.}
In a general AM protocol, deciding whether a
random string is accepting means deciding
$\exists y\,V(x,r,y)$; the response cannot be
enumerated in polynomial time. Substituting it
in the same covering formula yields
\[
 \exists(a_i)_i\ \forall r\ \exists(i,y)\quad
                   V'(x,r\mathbin\oplus a_i,y).
\]
The last existential block depends on $r$ and
cannot simply be moved in front of it. Listing
a response for each possible $r$ would use an
exponential-size table. Guessing a polynomial-size
circuit for those responses would require an
unproved bound on the complexity of a successful
Merlin strategy. Thus the argument has not
constructed the polynomial deterministic matrix
required for a two-block simulation.

\paragraph{Verification and final status.}
The root script exhaustively checks the translate
covering assertion on cubes through dimension
three, as a finite check on the quantifier
direction. The probabilistic calculation proves
the general covering lemma. The two logarithmic
restrictions are established, while the full
unconditional AM inclusion remains unresolved
by this attempt.

\subsection{Sources}
\begin{itemize}
\item[S1] ULAM.AI and the UnsolvedMath Contributors, supplied problems.json, record 3404 (OPG-454), OpenGarden collection. Catalogue identity and source transcription; CC BY 4.0.
\url{https://huggingface.co/datasets/ulamai/UnsolvedMath}.
\item[S2] Open Problem Garden, Unconditional derandomization of Arthur-Merlin games. Original problem and discussion; the mathematical formulation below is expanded from the supplied source record.
\url{http://www.openproblemgarden.org/op/unconditional_derandomization_of_mathcal_am}.

\end{itemize}
\end{document}
